% =============================================================================
% Master supplement for the pretrain-loss-weighting NN-XC manuscript.
% Build BEFORE main.tex — main.tex reads supplement.aux for xr-hyper cross-refs.
%
% Local build:
%   mkdir -p build_supp
%   pdflatex -output-directory=build_supp supplement.tex
%   (cd build_supp && BIBINPUTS=.. bibtex supplement)
%   pdflatex -output-directory=build_supp supplement.tex
%   pdflatex -output-directory=build_supp supplement.tex
%
% Overleaf:
%   - Set supplement.tex as the main document, compile.
%   - Then set main.tex as the main document; xr-hyper resolves cross-refs.
% =============================================================================
\documentclass[11pt,letterpaper]{article}

\usepackage[margin=1in]{geometry}
\usepackage{graphicx, booktabs, amsmath, amssymb}
\usepackage[hidelinks]{hyperref}
\usepackage{natbib, xcolor, microtype, caption}
\captionsetup{font=small, labelfont=bf}

% Renumber figures/tables as S1, S2, ... so cross-refs from main.tex
% render as "Fig. S<n>" / "Tab. S<n>" automatically.
\renewcommand{\thefigure}{S\arabic{figure}}
\renewcommand{\thetable}{S\arabic{table}}

\graphicspath{
  {step5_pretrain_origin_comparison/figures/}
  {step6_unweighted_analysis/figures/}
  {step6_integration_analysis/figures/}
  {step6_pretrain_origin_comparison/figures/}
  {step56_combined_synthesis/figures/}
  {../step5_pretrain_origin_comparison/figures/}
  {../step6_unweighted_analysis/figures/}
  {../step6_integration_analysis/figures/}
  {../step6_pretrain_origin_comparison/figures/}
  {../step56_combined_synthesis/figures/}
  {../../notebooks/checkpoints_step5/unweighted/figures/}
  {../../notebooks/checkpoints_step5/integration/figures/}
  {../../notebooks/checkpoints_step6/unweighted/figures/}
  {../../notebooks/checkpoints_step6/integration/figures/}
  {../}
}

\title{Supplementary Material: Pretrain-Loss-Weighting Choice for
       Neural-Network Density Functionals}
\author{Alec P.\ Wills}
\date{\today}

\begin{document}
\maketitle

\tableofcontents
\bigskip\hrule\bigskip

\section{Pretraining diagnostics}
\label{S:sec:pretraining}

This section documents the pretraining phase across all four workflow
combinations (step 5 unweighted, step 5 integration, step 6 unweighted,
step 6 integration). Headline numbers are reproduced in the main paper;
this supplement adds the per-architecture and diagnostic depth.

\subsection{Per-architecture $F_x$ and $F_c$ pretrain MSE}
\label{S:sec:pretrain_per_arch}

\begin{table}[h]
\centering
\caption{Step-5 per-architecture pretrain final-step MSE for $F_x$ and
$F_c$, both origins. Ratios $>$ 1 indicate integration tighter.
Source: \texttt{notebooks/checkpoints\_step5/\{unweighted,integration\}/pretrain/<arch>/pretrain\_metadata.json}.}
\label{S:tab:pretrain_per_arch}
\small
\begin{tabular}{lrrrrrr}
\toprule
arch & $F_x$ uw & $F_x$ int & $F_x$ ratio & $F_c$ uw & $F_c$ int & $F_c$ ratio \\
\midrule
deep                & $2.88\!\times\!10^{-4}$ & $6.07\!\times\!10^{-5}$ & 4.74 & $1.15\!\times\!10^{-4}$ & $7.03\!\times\!10^{-5}$ & 1.63 \\
deep\_attn          & $2.83\!\times\!10^{-4}$ & $5.97\!\times\!10^{-5}$ & 4.74 & $9.99\!\times\!10^{-5}$ & $7.66\!\times\!10^{-5}$ & 1.30 \\
deep\_dm            & $2.30\!\times\!10^{-4}$ & $5.65\!\times\!10^{-5}$ & 4.06 & $2.95\!\times\!10^{-5}$ & $4.87\!\times\!10^{-5}$ & 0.61 \\
deep\_dm\_attn      & $2.18\!\times\!10^{-4}$ & $5.41\!\times\!10^{-5}$ & 4.03 & $1.87\!\times\!10^{-5}$ & $2.21\!\times\!10^{-5}$ & 0.85 \\
deep\_cusp          & $5.79\!\times\!10^{-5}$ & $1.52\!\times\!10^{-5}$ & 3.82 & $3.79\!\times\!10^{-5}$ & $2.47\!\times\!10^{-5}$ & 1.53 \\
deep\_cusp\_attn    & $4.56\!\times\!10^{-5}$ & $1.01\!\times\!10^{-5}$ & 4.51 & $1.92\!\times\!10^{-5}$ & $2.18\!\times\!10^{-5}$ & 0.88 \\
deep\_combined      & $3.93\!\times\!10^{-5}$ & $1.33\!\times\!10^{-5}$ & 2.95 & $4.47\!\times\!10^{-5}$ & $8.66\!\times\!10^{-5}$ & 0.52 \\
deep\_combined\_attn& $3.51\!\times\!10^{-5}$ & $1.17\!\times\!10^{-5}$ & 3.01 & $1.95\!\times\!10^{-5}$ & $2.22\!\times\!10^{-5}$ & 0.88 \\
\midrule
mean (per-arch ratio) & & & 3.98 & & & 1.02 \\
ratio of population means & & & 4.25 & & & 1.03 \\
\bottomrule
\end{tabular}
\end{table}

\paragraph{$F_c$/$F_x$ asymmetry.} While integration tightens $F_x$ on every
architecture (8/8 step-5 + 2/2 step-6 = 10/10), $F_c$ is mixed: 5 of 8
architectures (\texttt{deep\_dm}, \texttt{deep\_dm\_attn},
\texttt{deep\_cusp\_attn}, \texttt{deep\_combined},
\texttt{deep\_combined\_attn}) are \emph{worse} under integration on $F_c$.
The 3 that benefit (\texttt{deep}, \texttt{deep\_attn}, \texttt{deep\_cusp})
include the two architectures with the largest absolute unweighted $F_c$
baselines ($1.15\!\times\!10^{-4}$ and $9.99\!\times\!10^{-5}$); their
ratios (1.63 and 1.30 in favour of integration) dominate the population mean.
Likely causes: (a) PW92 $\varepsilon_c$ is roughly an order of magnitude
smaller in magnitude than $\varepsilon_x^{\mathrm{LDA}}$ at typical bonding
densities, narrowing the $F_c$-weighted residual's dynamic range, and (b)
PBE's $F_c$ functional form is closer to the LDA limit than $F_x$ is in the
GGA regime, leaving less ``room to move'' under any re-weighting.

\subsection{Pretrain loss curves and parity plots}

\begin{figure}[h!]
\centering
\includegraphics[width=0.92\textwidth]{pretrain_losses.png}
\caption{Step-5 pretrain loss curves per architecture, both origins
(in-notebook output). $x$-axis: optimizer step.}
\label{S:fig:pretrain_losses_step5}
\end{figure}

\begin{figure}[h!]
\centering
\includegraphics[width=0.92\textwidth]{pretrain_parity.png}
\caption{Step-5 pretrain parity plots per architecture
(predicted vs.\ analytic-PBE $F_x$).}
\label{S:fig:pretrain_parity_step5}
\end{figure}

\begin{figure}[h!]
\centering
\includegraphics[width=0.6\textwidth]{aux_components_deep_combined.png}
\caption{Auxiliary loss components on the \texttt{deep\_combined}
architecture, step 5.}
\label{S:fig:aux_components_step5}
\end{figure}

\section{Optimization step 5: H$_2$O training}
\label{S:sec:step5}

\subsection{Headline numbers}
\label{S:sec:step5_headline}

\begin{table}[h]
\centering
\caption{Step-5 paired headline (verified against
\texttt{combined\_headline.json} and per-arch
\texttt{pretrain\_metadata.json}). Ratio columns: integration~$\div$~unweighted
where ``tighter'' means smaller MSE / smaller AE-MAE.}
\label{S:tab:step5_headline}
\small
\begin{tabular}{lrrr}
\toprule
metric & unweighted & integration & ratio \\
\midrule
best AE-MAE on H$_2$O (kcal/mol) & 0.001402 & \textbf{0.000118} & 11.9$\times$ tighter \\
best AE spec & deep\_dm/C/fixed\_j\_3 & deep\_dm\_attn/C/fixed\_j\_3 & adds attention \\
best $\rho$-RMSE on H$_2$O (e/bohr$^3$) & $4.41\!\times\!10^{-4}$ & $\mathbf{3.72\!\times\!10^{-4}}$ & 1.19$\times$ tighter \\
pretrain $F_x$ mean (8 archs) & $1.50\!\times\!10^{-4}$ & $\mathbf{3.52\!\times\!10^{-5}}$ & 4.25$\times$ tighter \\
pretrain $F_c$ mean (8 archs) & $4.80\!\times\!10^{-5}$ & $4.66\!\times\!10^{-5}$ & 1.03$\times$ (tied) \\
pretrained-only NN AE on H$_2$O (kcal/mol) & 27.2 & \textbf{17.1} & 1.59$\times$ tighter (37.1\%) \\
random-init NN AE on H$_2$O (kcal/mol) & 434.6 & 434.6 & tied \\
reduction vs.\ PBE on H$_2$O ($\sim$7.03 kcal/mol) & 5,014$\times$ & 59,500$\times$ & --- \\
\bottomrule
\end{tabular}
\end{table}

\subsection{Post-processing figures}

\begin{figure}[h!]
\centering
\includegraphics[width=\textwidth]{figs5_1_pretrain_loss.png}
\caption{Pretrain final $F_x$ (a) and $F_c$ (b) MSE per architecture, paired
bars (unweighted blue / integration green); see Tab.~\ref{S:tab:pretrain_per_arch}
for the numerics.}
\label{S:fig:s5_pretrain_bars}
\end{figure}

\begin{figure}[h!]
\centering
\includegraphics[width=\textwidth]{figs5_2_baseline_reduction.png}
\caption{H$_2$O AE-error baseline reduction. Trained-NN AE for all 72 specs
per origin; horizontal references are random-NN ($\sim$434~kcal/mol),
pretrained-only NN (27.2 vs.\ 17.1 kcal/mol), PBE-on-H$_2$O
($\sim$7.03~kcal/mol; this run's def2-svp / grid\_level=1 H$_2$O
configuration; \emph{not} a Karton 2011 literature value), and chemical
accuracy 1~kcal/mol \citep{Pople1999}.}
\label{S:fig:s5_baseline_reduction}
\end{figure}

\begin{figure}[h!]
\centering
\includegraphics[width=\textwidth]{figs5_3_pareto_density_vs_AE.png}
\caption{Density-vs-energy plane. Each point is one spec; colour = origin;
marker = architecture family. Integration green sits below-left of unweighted
blue --- a Pareto improvement on both axes simultaneously for most
architectures. Both clusters share the same density-RMSE floor at
$\sim 3.7\!\times\!10^{-4}$~e/bohr$^3$ (set by network expressive capacity
and quadrature-grid resolution, not by pretrain origin).}
\label{S:fig:s5_pareto}
\end{figure}

\begin{figure}[h!]
\centering
\includegraphics[width=\textwidth]{figs5_4_arch_landscape.png}
\caption{Full AE-MAE landscape: 8 archs $\times$ 9 (loss $\times$ solver) cells
per origin, log colour scale. Integration's lower-right (descriptor-rich
attention archs $\times$ \texttt{C\_atomization\_plus\_grid} loss) reaches
deep green (sub-millihartree), unmatched in the unweighted heatmap.}
\label{S:fig:s5_landscape}
\end{figure}

\begin{figure}[h!]
\centering
\includegraphics[width=\textwidth]{figs5_5_transfer.png}
\caption{Transfer-mol AE per architecture and origin on the primary transfer
set $\{$CH$_4$, H$_2$, OH$\}$. PBE references are hardcoded for our specific
def2-svp / grid\_level=1 configuration: CH$_4$ 0.91, H$_2$ 7.31, OH 1.84
kcal/mol. These are \emph{not} W4-11 literature values ---
\citet{KartonDaonMartin2011} tabulates CCSD(T)-quality reference atomization
energies, not PBE errors. No spec on any architecture or origin reaches
chemical accuracy on any of CH$_4$, H$_2$, OH: H$_2$O-only training is
insufficient for the closed-shell light-hydride class. OH is additionally
open-shell ($^2\Pi$, $S{=}1/2$); the pretrain atom set $\{$H, He, O, N$\}$
includes open-shell atoms, so the OH mismatch is on the molecular
spin-density profile, not on atomic open-shell character.}
\label{S:fig:s5_transfer}
\end{figure}

\subsection{In-notebook diagnostics}

The step-5 notebook (\texttt{notebooks/\_build\_step5\_notebook.py}) produces
the following diagnostics directly. They cover training-phase loss curves,
auxiliary-loss component breakdowns
(\texttt{aux\_components\_deep\_combined.png}, included in
\S\ref{S:sec:pretrain_per_arch} above), an exploratory loss-balancing pilot
(\texttt{balancing\_aux\_comparison.png},
\texttt{bal\_loss\_plot\_vxc.png} --- the latter motivated step 6's
$V_{xc}$-aware losses but is not part of step 5's production
$\{A, B, C\}$ loss family), GradNorm task-weight evolution, SCF convergence
aggregates, density-matrix and grid-density diagnostics, $F_x(s)$ drift on
the CH$_4$ molecular grid, transfer aggregate and per-mol breakdowns, and a
feature-impact analysis.

\begin{figure}[h!]\centering\includegraphics[width=0.92\textwidth]{training_losses.png}
\caption{Training-phase loss curves per architecture, loss, and solver.}
\label{S:fig:s5_training_losses}\end{figure}

\begin{figure}[h!]\centering\includegraphics[width=0.92\textwidth]{balancing_aux_comparison.png}
\caption{Loss-balancing comparison aux plot.}
\label{S:fig:s5_balancing_aux}\end{figure}

\begin{figure}[h!]\centering\includegraphics[width=0.6\textwidth]{bal_loss_plot_vxc.png}
\caption{Loss-balancing experiment loss curves on \texttt{deep\_combined}
(in-notebook). Exploratory pilot that motivated step-6 $V_{xc}$-aware loss
families; step 5's actual production losses $\{A, B, C\}$ are
exchange-and-density only.}
\label{S:fig:s5_bal_vxc}\end{figure}

\begin{figure}[h!]\centering\includegraphics[width=0.92\textwidth]{gradnorm_weight_evolution.png}
\caption{GradNorm task-weight evolution \citep{ChenGradNorm2018}.}
\label{S:fig:s5_gradnorm}\end{figure}

\begin{figure}[h!]\centering\includegraphics[width=0.92\textwidth]{scf_comparison_ae.png}
\caption{SCF-mode AE comparison.}
\label{S:fig:s5_scf_ae}\end{figure}

\begin{figure}[h!]\centering\includegraphics[width=0.6\textwidth]{scf_convergence.png}
\caption{SCF convergence aggregate (step 5).}
\label{S:fig:s5_scf_conv}\end{figure}

\begin{figure}[h!]\centering\includegraphics[width=0.92\textwidth]{dm_heatmaps_scf.png}
\caption{DM heatmaps after final SCF.}
\label{S:fig:s5_dm_heatmaps}\end{figure}

\begin{figure}[h!]\centering\includegraphics[width=0.92\textwidth]{grid_density_scf.png}
\caption{Grid-density profile after final SCF.}
\label{S:fig:s5_grid_density}\end{figure}

\begin{figure}[h!]\centering\includegraphics[width=0.6\textwidth]{fx_drift_ch4.png}
\caption{$F_x(s)$ drift sampled on the CH$_4$ molecular grid.}
\label{S:fig:s5_fx_drift_ch4}\end{figure}

\begin{figure}[h!]\centering\includegraphics[width=0.92\textwidth]{transfer_aggregate.png}
\caption{Transfer aggregate plot.}
\label{S:fig:s5_transfer_aggregate}\end{figure}

\begin{figure}[h!]\centering\includegraphics[width=0.6\textwidth]{transfer_CH4.png}
\caption{CH$_4$ transfer detail.}
\label{S:fig:s5_transfer_ch4}\end{figure}

\begin{figure}[h!]\centering\includegraphics[width=0.6\textwidth]{transfer_H2.png}
\caption{H$_2$ transfer detail.}
\label{S:fig:s5_transfer_h2}\end{figure}

\begin{figure}[h!]\centering\includegraphics[width=0.6\textwidth]{transfer_OH.png}
\caption{OH transfer detail.}
\label{S:fig:s5_transfer_oh}\end{figure}

\begin{figure}[h!]\centering\includegraphics[width=0.92\textwidth]{feature_impact_scf.png}
\caption{Feature-impact analysis (in-notebook).}
\label{S:fig:s5_feature_impact}\end{figure}

\section{Optimization step 6: H$_2$O+C$_2$H$_2$ training, paired origins}
\label{S:sec:step6}

\subsection{Step-6 paired headline}
\label{S:sec:step6_headline}

\begin{table}[h]
\centering
\caption{Step-6 paired comparison verified against
\texttt{step6\_pretrain\_origin\_comparison/headline\_diff.json}. Ratio
column: unweighted~$\div$~integration where applicable.}
\label{S:tab:step6_paired}
\small
\begin{tabular}{lrrr}
\toprule
metric & unweighted & integration & ratio (uw/int) \\
\midrule
\# specs & 90 & 90 & --- \\
best AE-MAE (kcal/mol) & 0.00793 & \textbf{0.00277} & 2.86$\times$ tighter \\
best AE spec & dca/L2/oneshot & dca/L5/full\_3 & --- \\
best $\rho$-RMSE (e/bohr$^3$) & $1.32\!\times\!10^{-4}$ & $\mathbf{1.23\!\times\!10^{-4}}$ & 1.07$\times$ tighter \\
random-init NN AE (kcal/mol) & 461.5 & 461.5 & tied \\
pretrained-only NN AE (kcal/mol) & 121.4 & \textbf{119.0} & 1.02$\times$ tighter \\
PBE AE-MAE on H$_2$O+C$_2$H$_2$ (kcal/mol) & 8.17 & 8.17 & tied (def2-svp/lvl1) \\
\midrule
g3 MAE \texttt{L1\_B} & 0.344 & 0.376 & 0.91$\times$ \\
g3 MAE \texttt{L2\_C\_anchor} & 0.337 & \textbf{0.099} & 3.39$\times$ tighter \\
g3 MAE \texttt{L3\_balanced\_vxc} & 1.616 & 2.418 & 0.67$\times$ \\
g3 MAE \texttt{L4\_balanced\_vxc\_anchor} & 2.150 & 2.747 & 0.78$\times$ \\
g3 MAE \texttt{L5\_gradnorm\_vxc} & 0.426 & 0.520 & 0.82$\times$ \\
\bottomrule
\end{tabular}
\end{table}

\paragraph{Per-mol PBE references.} 7.03~kcal/mol on H$_2$O alone,
9.89~kcal/mol on C$_2$H$_2$ alone (mean magnitude over the training set:
8.17~kcal/mol).

\clearpage
\subsection{Step-6 post-processing figures (paired by name)}

\begin{figure}[h!]\centering\includegraphics[width=\textwidth]{step6_unweighted_analysis/figures/fig1_baseline_reduction.png}
\caption{Baseline AE reduction ladder, step-6 unweighted origin.}
\label{S:fig:s6u_fig1}\end{figure}

\begin{figure}[h!]\centering\includegraphics[width=\textwidth]{step6_integration_analysis/figures/fig1_baseline_reduction.png}
\caption{Baseline AE reduction ladder, step-6 integration origin.}
\label{S:fig:s6i_fig1}\end{figure}

\begin{figure}[h!]\centering\includegraphics[width=\textwidth]{step6_unweighted_analysis/figures/fig1b_baseline_reduction_transfer.png}
\caption{Same ladder evaluated on the two transfer sets, step-6 unweighted.
Top row = primary $\{$CH$_4$, H$_2$, OH$\}$; bottom row = secondary
$\{$CO$_2$, HF, NH$_2$, NH$_3\}$.}
\label{S:fig:s6u_fig1b}\end{figure}

\begin{figure}[h!]\centering\includegraphics[width=\textwidth]{step6_integration_analysis/figures/fig1b_baseline_reduction_transfer.png}
\caption{Transfer ladder, step-6 integration.}
\label{S:fig:s6i_fig1b}\end{figure}

\begin{figure}[h!]\centering\includegraphics[width=\textwidth]{step6_unweighted_analysis/figures/fig2_density_vs_energy_tradeoff.png}
\caption{Density-vs-energy tradeoff, step-6 unweighted.}
\label{S:fig:s6u_fig2}\end{figure}

\begin{figure}[h!]\centering\includegraphics[width=\textwidth]{step6_integration_analysis/figures/fig2_density_vs_energy_tradeoff.png}
\caption{Density-vs-energy tradeoff, step-6 integration. The L3/L4 cluster
sits to the right (worse AE) but downward (tighter $\rho$-RMSE), exhibiting
the classic Burke--Perdew energy-vs-density tradeoff (paraphrased in
\citealp{Kepp2017}).}
\label{S:fig:s6i_fig2}\end{figure}

\begin{figure}[h!]\centering\includegraphics[width=\textwidth]{step6_unweighted_analysis/figures/fig3_in_dist_vs_transfer.png}
\caption{In-distribution vs.\ transfer-mol AE per spec, step-6 unweighted.}
\label{S:fig:s6u_fig3}\end{figure}

\begin{figure}[h!]\centering\includegraphics[width=\textwidth]{step6_integration_analysis/figures/fig3_in_dist_vs_transfer.png}
\caption{In-distribution vs.\ transfer-mol AE, step-6 integration.}
\label{S:fig:s6i_fig3}\end{figure}

\begin{figure}[h!]\centering\includegraphics[width=\textwidth]{step6_unweighted_analysis/figures/fig4_loss_strategy_heatmap.png}
\caption{Loss-strategy AE-MAE heatmap, step-6 unweighted. Rows = loss family;
columns = the 18 (group, arch, solver) cells on training mols (H$_2$O+C$_2$H$_2$).
Cell value = mean $|\!$AE error$|\!$; colour = $\log_{10}$ of that mean
(\texttt{notebooks/analysis/comparison\_lib.py:636--674}).}
\label{S:fig:s6u_fig4}\end{figure}

\begin{figure}[h!]\centering\includegraphics[width=\textwidth]{step6_integration_analysis/figures/fig4_loss_strategy_heatmap.png}
\caption{Loss-strategy heatmap, step-6 integration.}
\label{S:fig:s6i_fig4}\end{figure}

\begin{figure}[h!]\centering\includegraphics[width=\textwidth]{step6_unweighted_analysis/figures/fig5_arch_comparison.png}
\caption{Architecture comparison, step-6 unweighted. Bars (panel a, ``trained
mols'') plot the mean $|\!$AE error$|\!$ over all 9 (group, solver) cells per
(loss, arch). Group-3 medians (deep vs.\ attn): L1\_B 0.65/0.04,
L2\_C\_anchor 0.65/0.01, L3 2.51/1.23, L4 2.58/1.53, L5 0.83/0.02
(recomputed from \texttt{eval\_df.csv}).}
\label{S:fig:s6u_fig5}\end{figure}

\begin{figure}[h!]\centering\includegraphics[width=\textwidth]{step6_integration_analysis/figures/fig5_arch_comparison.png}
\caption{Architecture comparison, step-6 integration. Group-3 medians
(deep vs.\ attn): L1\_B 0.74/0.01, L2 0.15/0.02, L3 0.58/2.25,
L4 4.74/0.65, L5 1.03/0.003.}
\label{S:fig:s6i_fig5}\end{figure}

\clearpage
\subsection{Step-6 in-notebook diagnostics}

The step-6 notebook (\texttt{notebooks/\_build\_step6\_notebook.py})
produces these diagnostics during the run, independent of the
post-processing analysis. They cover pretrain dynamics, anchor and
balancing-loss sensitivity, $F_x(s)$ profiles on probe-molecule grids
(panel B = CH$_4$ + C$_2$H$_2$; panel C = held-out C$_2$H$_4$),
per-group training-loss curves, SCF convergence, transfer-set MAE
breakdowns, and a $V_{xc}$-loss efficacy diagnostic. Captions describe
each plot's semantics.

\begin{figure}[h!]\centering\includegraphics[width=0.92\textwidth]{../../notebooks/checkpoints_step6/unweighted/figures/pretrain_losses.png}
\caption{Pretrain loss curves, step-6 unweighted.}
\label{S:fig:s6u_pretrain_losses}\end{figure}

\begin{figure}[h!]\centering\includegraphics[width=0.92\textwidth]{../../notebooks/checkpoints_step6/integration/figures/pretrain_losses.png}
\caption{Pretrain loss curves, step-6 integration.}
\label{S:fig:s6i_pretrain_losses}\end{figure}

\begin{figure}[h!]\centering\includegraphics[width=0.92\textwidth]{../../notebooks/checkpoints_step6/unweighted/figures/anchor_effect.png}
\caption{Anchor-loss effect on AE/density tradeoff, step-6 unweighted.}
\label{S:fig:s6u_anchor}\end{figure}

\begin{figure}[h!]\centering\includegraphics[width=0.92\textwidth]{../../notebooks/checkpoints_step6/integration/figures/anchor_effect.png}
\caption{Anchor-loss effect, step-6 integration.}
\label{S:fig:s6i_anchor}\end{figure}

\begin{figure}[h!]\centering\includegraphics[width=0.92\textwidth]{../../notebooks/checkpoints_step6/unweighted/figures/balancing_effect.png}
\caption{Balancing-loss effect, step-6 unweighted.}
\label{S:fig:s6u_balancing}\end{figure}

\begin{figure}[h!]\centering\includegraphics[width=0.92\textwidth]{../../notebooks/checkpoints_step6/integration/figures/balancing_effect.png}
\caption{Balancing-loss effect, step-6 integration.}
\label{S:fig:s6i_balancing}\end{figure}

\begin{figure}[h!]\centering\includegraphics[width=0.92\textwidth]{../../notebooks/checkpoints_step6/unweighted/figures/fx_drift_panel_B.png}
\caption{$F_x(s)$ on (CH$_4$, C$_2$H$_2$) probe grids (panel B),
step-6 unweighted. All 90 trained models overlaid.}
\label{S:fig:s6u_fx_drift_B}\end{figure}

\begin{figure}[h!]\centering\includegraphics[width=0.92\textwidth]{../../notebooks/checkpoints_step6/integration/figures/fx_drift_panel_B.png}
\caption{$F_x(s)$ on (CH$_4$, C$_2$H$_2$) probe grids (panel B),
step-6 integration.}
\label{S:fig:s6i_fx_drift_B}\end{figure}

\begin{figure}[h!]\centering\includegraphics[width=0.92\textwidth]{../../notebooks/checkpoints_step6/unweighted/figures/fx_drift_panel_C.png}
\caption{$F_x(s)$ on the C$_2$H$_4$ held-out probe (panel C),
step-6 unweighted. C$_2$H$_4$ is in neither training nor transfer set.}
\label{S:fig:s6u_fx_drift_C}\end{figure}

\begin{figure}[h!]\centering\includegraphics[width=0.92\textwidth]{../../notebooks/checkpoints_step6/integration/figures/fx_drift_panel_C.png}
\caption{$F_x(s)$ on C$_2$H$_4$ held-out probe (panel C), step-6 integration.}
\label{S:fig:s6i_fx_drift_C}\end{figure}

\clearpage

\begin{figure}[h!]\centering\includegraphics[width=0.92\textwidth]{../../notebooks/checkpoints_step6/unweighted/figures/loss_curves_group1_h2o_short.png}
\caption{Loss curves, group~1 (H$_2$O, short schedule), step-6 unweighted.}
\label{S:fig:s6u_loss_g1}\end{figure}

\begin{figure}[h!]\centering\includegraphics[width=0.92\textwidth]{../../notebooks/checkpoints_step6/unweighted/figures/loss_curves_group2_h2o_c2h2_short.png}
\caption{Loss curves, group~2 (H$_2$O + C$_2$H$_2$, short), step-6 unweighted.}
\label{S:fig:s6u_loss_g2}\end{figure}

\begin{figure}[h!]\centering\includegraphics[width=0.92\textwidth]{../../notebooks/checkpoints_step6/unweighted/figures/loss_curves_group3_h2o_c2h2_long.png}
\caption{Loss curves, group~3 (H$_2$O + C$_2$H$_2$, long), step-6 unweighted.}
\label{S:fig:s6u_loss_g3}\end{figure}

\begin{figure}[h!]\centering\includegraphics[width=0.92\textwidth]{../../notebooks/checkpoints_step6/integration/figures/loss_curves_group1_h2o_short.png}
\caption{Loss curves, group~1, step-6 integration.}
\label{S:fig:s6i_loss_g1}\end{figure}

\begin{figure}[h!]\centering\includegraphics[width=0.92\textwidth]{../../notebooks/checkpoints_step6/integration/figures/loss_curves_group2_h2o_c2h2_short.png}
\caption{Loss curves, group~2, step-6 integration.}
\label{S:fig:s6i_loss_g2}\end{figure}

\begin{figure}[h!]\centering\includegraphics[width=0.92\textwidth]{../../notebooks/checkpoints_step6/integration/figures/loss_curves_group3_h2o_c2h2_long.png}
\caption{Loss curves, group~3, step-6 integration.}
\label{S:fig:s6i_loss_g3}\end{figure}

\clearpage

\begin{figure}[h!]\centering\includegraphics[width=0.6\textwidth]{../../notebooks/checkpoints_step6/unweighted/figures/scf_convergence.png}
\caption{SCF-mode convergence, step-6 unweighted.}
\label{S:fig:s6u_scf}\end{figure}

\begin{figure}[h!]\centering\includegraphics[width=0.6\textwidth]{../../notebooks/checkpoints_step6/integration/figures/scf_convergence.png}
\caption{SCF-mode convergence, step-6 integration.}
\label{S:fig:s6i_scf}\end{figure}

\begin{figure}[h!]\centering\includegraphics[width=0.92\textwidth]{../../notebooks/checkpoints_step6/unweighted/figures/transfer_primary_mae.png}
\caption{Primary transfer-set MAE per spec, step-6 unweighted.}
\label{S:fig:s6u_transfer_primary}\end{figure}

\begin{figure}[h!]\centering\includegraphics[width=0.92\textwidth]{../../notebooks/checkpoints_step6/integration/figures/transfer_primary_mae.png}
\caption{Primary transfer-set MAE, step-6 integration.}
\label{S:fig:s6i_transfer_primary}\end{figure}

\begin{figure}[h!]\centering\includegraphics[width=0.92\textwidth]{../../notebooks/checkpoints_step6/unweighted/figures/transfer_secondary_mae.png}
\caption{Secondary transfer-set MAE per spec, step-6 unweighted.}
\label{S:fig:s6u_transfer_secondary}\end{figure}

\begin{figure}[h!]\centering\includegraphics[width=0.92\textwidth]{../../notebooks/checkpoints_step6/integration/figures/transfer_secondary_mae.png}
\caption{Secondary transfer-set MAE, step-6 integration.}
\label{S:fig:s6i_transfer_secondary}\end{figure}

\begin{figure}[h!]\centering\includegraphics[width=0.92\textwidth]{../../notebooks/checkpoints_step6/unweighted/figures/vxc_efficacy.png}
\caption{$V_{xc}$ efficacy diagnostic, step-6 unweighted.}
\label{S:fig:s6u_vxc}\end{figure}

\begin{figure}[h!]\centering\includegraphics[width=0.92\textwidth]{../../notebooks/checkpoints_step6/integration/figures/vxc_efficacy.png}
\caption{$V_{xc}$ efficacy diagnostic, step-6 integration.}
\label{S:fig:s6i_vxc}\end{figure}

\bibliographystyle{plainnat}
\bibliography{references}

\end{document}
